\chapter{绪论 讲座内容多维概括}

\section{课程背景与总体框架}

学科技术前沿讲座课程旨在为计算机学院本科生提供了解学科最新发展动态、拓展学术视野、培养科研兴趣的平台。2026年春季学期，来自北航计算机学院及相关单位的八位专家学者围绕计算机科学的哲学基础、科研方法、人工智能（含大模型、具身智能）、虚拟现实、高性能计算、视觉计算、新型网络架构等主题做了系列报告。本报告是对这些讲座内容的系统总结与个人思考的集成\cite{li2026ai}。

\section{学科通论与科研方法论}

李波老师的开篇讲座《计算机科学的科研工作和论文写作》从哲学高度切入\ref{李波}，探讨是否存在计算机科学这一根本问题，辨析广义计算机科学（与数学、物理、化学、社会科学交叉）与狭义计算机科学（聚焦计算机与计算本身）的矛盾。他指出当前我国计算机领域普遍存在重应用轻核心的现象，狭义计算机科学的研究人员过剩，而能够突破核心技术的人才匮乏。李老师认为科研的本质是“Search and Search Again”，创新点的提出是最大难点。在论文写作方面，他强调透彻的背景分析（Summary, Survey, Motivation）、翔实的数学模型（数学公式往往是本质性的）以及有代表性的实验数据的重要性。这些观点为我们后续理解各具体领域的技术前沿提供了方法论指导。

\section{人工智能 从大模型到具身智能}

\subsection{人工智能发展状况与应用特点}

李波老师在另一场讲座《面向智能化发展的科技创新与人才培养》中系统梳理了人工智能从机械革命、电力革命、信息革命到当前智能化革命的演进脉络。他特别指出尽管人工智能潜力巨大，实际项目中人工智能的失败率高达百分之八十以上，主要原因在于数据问题、方法选择错误以及未能与具体业务场景深度结合。他提出二八现象，即通用大模型可能适合百分之八十的大众化应用，而百分之二十的复杂特殊问题需要专精特新的垂直领域模型。他还深入分析了人工智能面临的常识推理、不确定性推理等长期难题，以及深度学习作为静态数据拟合方法在安全攸关任务中的局限性，引导我们理性看待人工智能的能力与边界。

\subsection{大模型技术与DeepSeek案例}

李舟军教授的讲座《DeepSeek那些事儿》以现象级产品DeepSeek为案例，深入讲解大模型技术原理。他回顾了从专家系统到机器学习再到预训练加后训练范式的人工智能发展简史，重点剖析DeepSeek的技术创新，包括混合专家架构带来的训练与推理成本优势、GRPO强化学习方法赋予模型的强推理能力，以及全面开源策略带来的鲸鱼效应。李教授还结合自身团队的工作展示大模型在代码生成、社交网络内容理解、智能运维、网络威胁情报分析等垂直领域的落地应用新范式，让我们看到大模型从通用聊天走向行业专家的巨大潜力。

\subsection{具身智能与智能无人机}

刘偲老师的讲座《具身智能与智能无人机》将人工智能延伸到物理世界。她详细阐述具身智能的三大核心要素，包括智能体（由大语言模型或视觉语言模型驱动的大脑）、本体（人形机器人、无人机、机械臂等身体）以及数据与进化架构（仿真环境如Omniverse）。她介绍国内外在具身智能领域的战略布局，并展示团队在无人机视觉语言动作模型、长距离规划导航以及多智能体协同等方面的前沿研究成果，生动诠释人工智能如何从虚拟大脑进化为能与物理世界交互的智能实体。

\section{虚拟现实与数字孪生}

周忠教授的讲座《虚拟现实技术前沿》系统介绍虚拟现实与增强现实技术的发展脉络、核心概念（沉浸感、交互性、构想性以及虚拟现实2.0提出的智能化、互通性、迭代性）与最新应用。他从太空算力、机器人谈到虚拟现实与增强现实在民生、工业、航天等领域的应用实例，深入介绍数字孪生技术及其在国防、地球科学、智能制造中的颠覆性作用。周教授还展示北航在虚拟现实专业课程建设与跨学科大作业方面的成果，体现虚拟现实作为解决问题思路的广阔前景。

\section{多维度视觉计算}

黄迪教授的讲座《多维度视觉计算》聚焦计算机视觉这一人工智能的核心领域。他从视觉计算的三个维度（目标、行为、事件）出发，探讨核心问题即视觉表征，介绍从生物启发的传统方法到基于预训练大模型的现代方法演进。黄教授团队在目标检测、识别、跟踪、分割、弱监督学习以及视觉内容生成等方面取得系列前沿成果，展现视觉计算从感知到生成再到理解的全面进步。

\section{高性能计算}

李波老师的讲座《从天河超级计算机看高性能计算技术》带我们回顾计算能力的基石。他以蛋白质结构预测、核爆模拟等大科学工程为引，阐述高性能计算对于国家战略与科技发展的重要性。讲座详细介绍天河一号、二号以及Blue Gene/L、Cray RedStorm、地球模拟器等典型机型架构，深入分析高性能计算的核心技术挑战，包括如何实现高生产率、高效率以及控制功耗。他探讨当前主流的集群体系结构、异构计算、高速互连网络、并行文件系统等技术趋势，让我们深刻理解支撑上层所有智能应用的算力是如何炼成的。

\section{存算传融合的自主计算}

张珊教授的讲座《存算传融合的自主计算》从网络互联的基本问题出发，回顾电路交换与分组交换的权衡，以及互联网的命名、寻址、路由、拥塞控制等经典难题。她敏锐指出当前传送算力存储三大资源体系独立演化的局限性，各自为政导致性能瓶颈。为此她提出存算传融合的创新理念，展示在芯片层面、架构层面以及系统层面的融合尝试。张教授描绘自主计算的愿景，即网络与设备融合后能够根据用户需求和环境状态自主按需安全地构建黑盒式计算环境。她还介绍团队在面向高时效视频分析的传算资源拟合、推理任务批处理机制等方面的具体研究成果。

\section{小结}

以上八位老师的讲座从哲学方法论到基础理论，从核心算法到系统架构，从软件智能到物理实体，共同勾勒出一幅计算机科学前沿发展的全景图。这些内容既相互独立又内在关联，为我后续深入思考个人收获与未来方向提供了丰富的素材。